\section{The generic form of the realization theorems}
\label{sec:generic-lift}

The realization theorems of the preceding chapters are existence
statements. Tay's theorem (\cref{thm:tay-witness}) and the body-hinge
Tay--Whiteley theorem (\cref{thm:body-hinge-tay}) construct a single
independent (resp.\ infinitesimally rigid) realization exactly when
the tree-packing count holds, and the molecule rank formula
(\cref{cor:molecule-rank-formula}) computes the rank of the generic
bar-joint matroid rather than the rank of any one realization. The
classical statements are stronger: \emph{almost all} realizations
attain the generic rank. This chapter proves that stronger form,
following the coordinate approach of Jackson and
Jord\'an~\cite{jacksonJordan2009}, which avoids the
irreducible-variety argument of
Whiteley~\cite[Proposition~6]{whiteley1988} that the body-bar and
body-hinge chapters left aside.

The pattern is uniform across the framework classes, and is the one
already in place for bar-joint placements
(\cref{sec:bar-joint-3d-generic-placement}). A realization is
\emph{generic} when every family of rigidity-matrix rows that is
linearly independent at some realization is linearly independent at
it (\cref{def:generic-placement} is the bar-joint instance). Such a
realization is in particular generic in the sense of Jackson and
Jord\'an --- its rigidity matrix and all edge-induced submatrices
attain their maximum rank over all
realizations~\cite[\S5]{jacksonJordan2009} --- and for classes with
one rigidity row per edge the two conditions coincide. Every
existence theorem then upgrades at once: a generic realization has
rank at least that of any witness realization, so a rank attained
once is attained at every generic realization.

Two statements quantify how common generic realizations are. They
exist (\cref{lem:exists-generic-placement}); and, since the entries
of the rigidity matrix are polynomials in the realization parameters,
the non-generic realizations are confined to the zero set of a single
nonzero polynomial (\cref{lem:generic-placement-abundance}). Over
$\R$ the complement of that zero set is dense and open, which is the
precise content of ``almost all realizations are generic''; Jackson
and Jord\'an instead note that realizations whose coordinates are
algebraically independent over $\Q$ are
generic~\cite[\S5]{jacksonJordan2009}.

The molecular form comes first (\cref{sec:generic-lift-molecule}):
every generic bar-joint realization of a square $G^2$ realizes the
molecule rank formula, and in particular is infinitesimally rigid
whenever $5 \cdot G$ contains six edge-disjoint spanning trees. The
panel-and-hinge form (over the panel normals), the body-bar form
(over the bar endpoints), and the body-and-hinge form (over points
spanning the hinges) follow the same pattern over their own
realization parameters.

\subsection{Abundance of generic placements}
\label{sec:generic-lift-abundance}

The existence lemma (\cref{lem:exists-generic-placement})
interpolates along one line at a time. The abundance form
strengthens it: one polynomial witnesses genericity across the whole
placement space.

\begin{lemma}[Abundance of generic placements]
  \label{lem:generic-placement-abundance}
  \lean{SimpleGraph.exists_isGenericPlacement_abundance}
  \leanok
  \uses{def:generic-placement}
  Let $V$ be finite and $d$ a dimension. There is a nonzero real
  polynomial $P$ in the $d\,|V|$ coordinates of a placement such that
  every placement $p : V \to \R^d$ with $P(p) \ne 0$ is generic for
  row independence. In particular the non-generic placements are
  contained in the zero set of a nonzero polynomial.
\end{lemma}
\begin{proof}
  \leanok
  \uses{def:edgeSet-rowIndependent}
  For an edge subset $I \subseteq E(K_V)$ that is row-independent at
  some placement $q_I$, the rigidity rows indexed by $I$ are linear in
  the placement coordinates, so a maximal square minor of the row
  family, chosen nonsingular at $q_I$, is a nonzero polynomial $P_I$ in
  the placement coordinates whose non-vanishing at a placement is
  sufficient for $I$ to remain row-independent there. Let $P$ be the
  product of the $P_I$ over
  the finitely many such subsets $I$; a finite product of nonzero
  polynomials is nonzero. If $P(p) \ne 0$ then every factor
  $P_I(p) \ne 0$, so every edge subset row-independent at some
  placement is row-independent at $p$.
\end{proof}

\subsection{Every generic realization of a square}
\label{sec:generic-lift-molecule}

At a placement generic for row independence the realized rank of
every graph equals its generic rank
(\cref{lem:rigidityMap-rank-at-generic}), so the molecule rank
formula (\cref{cor:molecule-rank-formula}) --- a statement about the
generic matroid --- becomes a statement about every generic
realization of the square.

\begin{theorem}[The molecule rank formula at every generic placement]
  \label{thm:molecule-generic-rank}
  \lean{SimpleGraph.molecule_generic_rank}
  \leanok
  \uses{def:generic-placement, def:square-graph, def:D-deficiency,
    def:rigidityMap}
  Let $G$ be a simple graph of minimum degree at least two on a
  finite, nonempty vertex set $V$, and let $p : V \to \R^3$ be a
  placement generic for row independence. Then
  \[
    \operatorname{rank} R(G^2, p)
      \;=\; 3|V| - 6 - \operatorname{def}(\tilde G),
  \]
  where $\tilde G = 5 \cdot G$.
\end{theorem}
\begin{proof}
  \leanok
  \uses{cor:molecule-rank-formula, lem:rigidityMap-rank-at-generic}
  At a generic placement the rank of the rigidity matrix of $G^2$
  equals the generic rank $r(G^2)$
  (\cref{lem:rigidityMap-rank-at-generic}), and
  $r(G^2) = 3|V| - 6 - \operatorname{def}(\tilde G)$ is
  \cref{cor:molecule-rank-formula}.
\end{proof}

\begin{corollary}[Every generic realization of a rigid square]
  \label{cor:molecule-generic-rigid}
  \lean{SimpleGraph.molecule_generic_rigid}
  \leanok
  \uses{def:generic-placement, def:square-graph, def:D-deficiency,
    def:isInfinitesimallyRigid}
  Let $G$ be a simple graph of minimum degree at least two on a
  finite, nonempty vertex set $V$, and suppose
  $\operatorname{def}(\tilde G) = 0$ for $\tilde G = 5 \cdot G$. Then
  for every placement $p : V \to \R^3$ generic for row independence
  the framework $(G^2, p)$ is infinitesimally rigid.
\end{corollary}
\begin{proof}
  \leanok
  \uses{thm:molecule-generic-rank}
  By \cref{thm:molecule-generic-rank},
  $\operatorname{rank} R(G^2, p) = 3|V| - 6$, so by rank--nullity in
  the $3|V|$-dimensional space of motions
  $\dim \ker R(G^2, p) = 6$, the bound of
  \cref{def:isInfinitesimallyRigid}.
\end{proof}

Jackson and Jord\'an observed~\cite[p.~586]{jacksonJordan2009} that
the molecular conjecture, then open, would imply that every generic
bar-joint realization of $G^2$ in $\R^3$ is infinitesimally rigid
whenever $5 \cdot G$ contains six edge-disjoint spanning trees. The
conjecture is a theorem of this development
(\cref{thm:molecular-conjecture}), so the implication holds
unconditionally.

\begin{corollary}[Generic rigidity of squares from six spanning
  trees; Jackson--Jord\'an]
  \label{cor:molecule-generic-square-packing}
  \lean{SimpleGraph.molecule_generic_square_packing}
  \leanok
  \uses{def:generic-placement, def:square-graph,
    def:isInfinitesimallyRigid}
  Let $G$ be a simple graph on a finite, nonempty vertex set $V$, and
  suppose $5 \cdot G$ contains six edge-disjoint spanning trees. Then
  every bar-joint realization of $G^2$ in $\R^3$ at a placement
  generic for row independence is infinitesimally rigid.
\end{corollary}
\begin{proof}
  \leanok
  \uses{cor:molecule-generic-rigid, thm:def-eq-corank, def:matroid-MG,
    def:cut-edges-2ec}
  For $|V| = 1$ the square has no edges and every placement is
  infinitesimally rigid ($\dim \ker R(G^2, p) \le 3 \le 6$,
  independently of the packing hypothesis). For $|V| \ge 2$: each of
  the six spanning trees has exactly $|V| - 1$ edges and, being
  acyclic, is independent in the cycle matroid of $5 \cdot G$; the six
  trees being pairwise edge-disjoint, their union is an independent
  set of size $6(|V| - 1)$ in the sixfold cycle-matroid union, hence
  in $M(\tilde G)$ (\cref{def:matroid-MG}). Since
  $\operatorname{def}(\tilde G) \ge 0$ caps
  $\operatorname{rank} M(\tilde G) \le 6(|V| - 1)$, this independent
  set is a base, so $\operatorname{def}(\tilde G) = 0$ by
  \cref{thm:def-eq-corank}. A body-hinge-rigid ($0$-dof) graph is
  $2$-edge-connected (\cref{def:cut-edges-2ec}), so every vertex of
  $5 \cdot G$ has degree at least two, hence every vertex of $G$ has
  degree at least two (the shadowing carrier's degree matches $G$'s
  one-for-one); \cref{cor:molecule-generic-rigid} applies.
\end{proof}
